% config/environments.tex — Theorem family, examples, exercises, boxes
% -----------------------------------------------------------------------
% Requires (loaded in main.tex before this file):
%   amsmath, amsthm, thmtools, tcolorbox, enumitem, environ,
%   config/colors, config/commands
%
% Design language:
%   * Definition family : green edge, pale green wash  ("growing ideas")
%   * Theorem family    : navy edge, pale blue wash    ("core results")
%   * Remarks/strategies: slate frame, faint wash      ("guidance")
%   * Exercises         : slate frame on white         ("your turn")
%   * Challenges        : gold edge, pale gold wash    ("summit")
%   * Solutions         : blue-gray wash (instructor edition only)
% All boxes are BREAKABLE — they split cleanly across pages.

% =====================  tcolorbox base styles  =========================
\usepackage{thmtools}     % \declaretheorem (requires amsthm, already loaded)
\usepackage{tcolorbox}
\tcbuselibrary{skins,breakable,theorems}
\tcbset{
  curriculum/.style={
    enhanced, breakable,
    fonttitle=\sffamily\bfseries\small,
    coltitle=white,
    colback=white,
    boxrule=0.5pt, arc=1.6pt,
    left=9pt, right=9pt, top=7pt, bottom=7pt,
    before skip=13pt plus 2pt, after skip=13pt plus 2pt,
    toptitle=3.5pt, bottomtitle=3pt,
  },
  defbox/.style={curriculum, colframe=forest, colback=forestlight,
    leftrule=2.6pt, rightrule=0.5pt, toprule=0.5pt, bottomrule=0.5pt},
  thmbox/.style={curriculum, colframe=navy, colback=navylight,
    leftrule=2.6pt, rightrule=0.5pt, toprule=0.5pt, bottomrule=0.5pt},
  remarkbox/.style={curriculum, colframe=slate, colback=slatelight,
    boxrule=0.5pt, colbacktitle=slate},
  exbox/.style={curriculum, colframe=slate!75, colback=white,
    boxrule=0.6pt, colbacktitle=slate!75,
    titlerule=0.5pt, titlerule style={slate!45}},
  chalbox/.style={curriculum, colframe=gold, colback=goldlight,
    leftrule=2.6pt, rightrule=0.5pt, toprule=0.5pt, bottomrule=0.5pt},
  solbox/.style={curriculum, colframe=solutionblue!90, colback=solutionlight,
    boxrule=0.5pt, colbacktitle=solutionblue!90},
}

% =====================  Theorem machinery  =============================
% Shared counter across theorem/lemma/proposition/corollary; definitions
% and examples carry their own per-chapter counters.
\declaretheorem[style=plain, name=Theorem, parent=chapter]{theorem}
\declaretheorem[style=plain, name=Lemma, sibling=theorem]{lemma}
\declaretheorem[style=plain, name=Proposition, sibling=theorem]{proposition}
\declaretheorem[style=plain, name=Corollary, sibling=theorem]{corollary}
\declaretheorem[style=definition, name=Definition, parent=chapter]{definition}
\declaretheorem[style=definition, name=Example, parent=chapter]{example}
\declaretheorem[style=remark, name=Remark, numbered=unless unique,
  parent=chapter]{remark}
\declaretheorem[style=remark, name=Key Idea, parent=chapter]{keyidea}
\declaretheorem[style=remark, name=Strategy, parent=chapter]{strategy}
\declaretheorem[style=remark, name=Common Pitfall, parent=chapter]{pitfall}

% Wrap the named environments in their boxes. Proofs stay unboxed
% (amsthm), typeset in the standard way for maximal readability.
\tcolorboxenvironment{theorem}{thmbox}
\tcolorboxenvironment{lemma}{thmbox}
\tcolorboxenvironment{proposition}{thmbox}
\tcolorboxenvironment{corollary}{thmbox}
\tcolorboxenvironment{definition}{defbox}
\tcolorboxenvironment{example}{exbox}
\tcolorboxenvironment{remark}{remarkbox}
\tcolorboxenvironment{keyidea}{remarkbox}
\tcolorboxenvironment{strategy}{remarkbox}
\tcolorboxenvironment{pitfall}{%
  curriculum, colframe=crimson, colback=crimsonlight,
  leftrule=2.6pt, rightrule=0.5pt, toprule=0.5pt, bottomrule=0.5pt,
  colbacktitle=crimson}

% Worked solutions inside examples: proof-like block with custom head.
\newenvironment{exsolution}{\begin{proof}[\normalfont\sffamily\bfseries\color{navy}Solution.]}{\end{proof}}

% =====================  Exercises with difficulty  =====================
% \begin{exercise}{3}                     -> 3-star exercise
% \begin{exercise}[label=ex:mylabel]{3}   -> add \label for solutions
% Optional tcolorbox keys go in the bracketed argument.
\newtcolorbox[auto counter, number within=chapter]{exercise}[2][]{%
  exbox,
  title={Exercise \thetcbcounter\ \hfill \difficulty{#2}},
  breakable, enhanced,
  #1
}

% Challenge problems: gold boxes, separate counter. Stars are the
% white variant: gold stars would vanish on the gold title bar.
\newtcolorbox[auto counter, number within=chapter]{challenge}[2][]{%
  chalbox,
  title={Challenge \thetcbcounter\ \hfill \difficultywhite{#2}},
  breakable, enhanced,
  #1
}

% =====================  Hints & solutions  =============================
% Hint: compact slate box for inline or collected hints.
\newtcolorbox{hint}[1][]{%
  remarkbox, coltitle=white,
  left=7pt, right=7pt, top=5pt, bottom=5pt,
  before skip=9pt, after skip=9pt,
  title={Hint}, #1
}

% Solution: matches a labeled exercise; instructor edition only.
% Usage: \begin{solution}{ex:1-3} ... \end{solution}
\newtcolorbox{solution}[1]{%
  solbox,
  title={Solution to Exercise~\ref{#1}},
}

% =====================  Chapter furniture  =============================
% Learning objectives: opening box with checked items.
\newtcolorbox{objectives}[1][]{%
  thmbox, title={Learning Objectives}, #1
}
% Objective item with gold check bullet (use inside itemize).
\newcommand{\objective}{\item[\textcolor{gold}{$\checkmark$}]}

% Prerequisites: compact slate box.
\newtcolorbox{prerequisites}[1][]{%
  remarkbox, title={Prerequisites}, #1
}

% Notation table box.
\newtcolorbox{notation}[1][]{%
  remarkbox, title={Notation}, #1
}

% Chapter summary: end-of-chapter distilled results.
\newtcolorbox{chaptersummary}[1][]{%
  thmbox, title={Chapter Summary}, #1
}
% Summary item with navy triangle bullet (use inside itemize).
\newcommand{\summaryitem}{\item[\textcolor{navy}{\small$\blacktriangleright$}]}

% =====================  Answer key helper  =============================
% \answer{42} — quick answer line for answer keys.
\newcommand{\answer}[1]{\par\smallskip\noindent\labeltext{Answer:} #1\par}

% =====================  Instructor-only note  ==========================
% Visible only in the instructor edition; absorbed silently otherwise.
\ifstudent
  \NewEnviron{instructornote}{}%
\else
  \newtcolorbox{instructornote}[1][]{%
    remarkbox, colbacktitle=crimson, title={Instructor Note}, #1}
\fi
